\section{Experiments}
\label{sec:experiments}

\subsection{Setup}

We fine-tune Qwen-2.5-7B-Instruct on 8K multi-turn conversations from the UltraChat training split, a dataset of naturalistic dialogues across diverse topics. LoRA adapters (rank 8) are applied to all linear layers. The combined loss $\mathcal{L} = \mathcal{L}_\text{NTP} + \lambda_\text{GAR}\mathcal{L}_\text{GAR} + \lambda_\text{div}\mathcal{L}_\text{div}$ uses $\lambda_\text{GAR} = \lambda_\text{div} = 1.0$ with MSE target $\tau = 0.5$ (\S\ref{sec:loss}). We train for 3 epochs with AdamW (lr $= 2 \times 10^{-4}$, batch size 1, gradient accumulation 8, seed 42); each variant requires approximately 6 hours on a single NVIDIA RTX PRO 6000 Blackwell Server Edition GPU (102\,GB). Full hyperparameters are provided in Appendix~A.

To isolate the contribution of gating, we evaluate four configurations: (1)~\textbf{ThreadGate}: LoRA + per-layer learned gates + NTP + GAR loss; (2)~\textbf{Uniform Gate}: LoRA + a single shared gate across all layers + NTP + GAR loss; (3)~\textbf{LoRA-only}: LoRA + NTP + GAR loss, with gates fixed at 1.0 (no learned gating mechanism); (4)~\textbf{Base}: the pretrained model without any fine-tuning. Each configuration is evaluated under two teacher-forced protocols. \emph{IID evaluation} computes metrics on UltraChat text from the same distribution as training. \emph{Held-out evaluation} uses 50 non-overlapping conversations to test within-distribution generalization. We report three metrics: GAR-ours (attention to goal tokens at the last token of each assistant turn, as defined in \S\ref{sec:background}), GAR-Dongre \citep{dongre2026attention} (attention to goal tokens averaged across all response tokens, all layers, and all heads), and perplexity (PPL).

\subsection{Main Results}

Table~\ref{tab:four-way} presents the four-way ablation. The central finding is that all three fine-tuned variants converge to nearly identical performance across both evaluation settings. On held-out data, GAR-ours spans only 0.75 percentage points (pp) across ThreadGate (0.50), Uniform Gate (0.51), and LoRA-only (0.50), and PPL is identical at 2.78 for all three. This convergence replicates on IID evaluation, where GAR-ours is 0.49 for all three variants and PPL varies by at most 0.01. Neither per-layer gating (ThreadGate) nor shared gating (Uniform Gate) provides any measurable benefit over LoRA-only fine-tuning. These results hold without exception across every metric-split combination in Table~\ref{tab:four-way}.

The Base model confirms that LoRA fine-tuning is the operative variable. GAR-ours increases from 0.01 (Base) to approximately 0.49 (all fine-tuned variants) under teacher-forced evaluation, a roughly 45-fold increase. Since all three fine-tuned variants share the same GAR auxiliary loss ($\tau = 0.5$), this convergence near the target is expected; it confirms that the loss functions normally but does not by itself demonstrate benefit from the gating mechanism. PPL improves correspondingly from 4.23 to 2.63 on IID data, which confirms that NTP training functions normally across all configurations. The three-variant convergence holds across both metrics and both evaluation splits; overfitting cannot account for a pattern this consistent.

GAR-Dongre remains uniformly low across all conditions (0.00 for Base, 0.01--0.02 for fine-tuned variants). This is consistent with the observation of \citet{dongre2026attention} that the original GAR definition, which averages over all response tokens across all layers and heads, captures a persistently diffuse attention signal. The contrast with GAR-ours reflects measurement granularity: GAR-ours isolates the last token of each assistant turn, where learned attention steering concentrates its effect, while GAR-Dongre dilutes the signal across all positions.

The teacher-forced GAR-ours values of 0.49--0.51 stand in sharp contrast to autoregressive generation, where GAR-ours drops to 0.016 for ThreadGate. This 30$\times$ gap between teacher-forced and model-generated evaluation is the central empirical finding of this paper. Under teacher-forcing, the model conditions on UltraChat reference text as the prefix for each response; under autoregressive generation, the prefix consists of the model's own prior outputs. The attention patterns optimized for the former distribution do not transfer to the latter, and the gating mechanism does nothing to bridge this gap. We defer the mechanistic analysis of this distribution shift to \S\ref{sec:analysis}.

\begin{table*}[t]
\centering
\small
\caption{Four-way ablation on Qwen-2.5-7B-Instruct. IID evaluation uses teacher-forced UltraChat text; held-out uses 50 non-overlapping conversations. All three fine-tuned variants converge to within 0.75\,pp on GAR-ours and 0\,pp on PPL (held-out). GAR-O = GAR-ours; GAR-D = GAR-Dongre \citep{dongre2026attention}. No values are bolded because no configuration achieves a meaningful advantage.}
\label{tab:four-way}
\begin{tabular}{l ccc ccc}
\toprule
& \multicolumn{3}{c}{\textbf{IID Evaluation}} & \multicolumn{3}{c}{\textbf{Held-out Evaluation}} \\
\cmidrule(lr){2-4} \cmidrule(lr){5-7}
\textbf{Configuration} & GAR-O & GAR-D & PPL & GAR-O & GAR-D & PPL \\
\midrule
Base          & 0.01 & 0.00 & 4.23 & 0.01 & 0.00 & 4.34 \\
ThreadGate    & 0.49 & 0.02 & 2.63 & 0.50 & 0.02 & 2.78 \\
Uniform Gate  & 0.49 & 0.01 & 2.62 & 0.51 & 0.01 & 2.78 \\
LoRA-only     & 0.49 & 0.01 & 2.63 & 0.50 & 0.01 & 2.78 \\
\bottomrule
\end{tabular}
\end{table*}

\subsection{Cross-Model Validation}
\label{sec:cross-model}

To verify that the convergence pattern and the 30$\times$ training-inference gap are not specific to the Qwen architecture, we are conducting parallel experiments on Llama-3.1-8B-Instruct. Llama-3.1 differs from Qwen-2.5 in several architectural choices, including grouped-query attention with 8 KV heads and a distinct BPE tokenizer. These differences provide a meaningful test of generalization across model families. Results will be reported in the final version of this paper.